% ===================================================================
%  GAMBAR No. 7 — Desa ing mangsa adhem lan omah tani ing mangsa semi.
%
%  Traduction, faite en 2026, en javanais d'aujourd'hui, sur l'ido de
%  texto/io. Regles de la colonne : voir 10-gambar-01.tex. Deux
%  scenes, deux numerotations. Ouverture de la DEUXIEME SERIE du
%  volume : « SERI KAPINDHO ».
%
%  UN BLOC A DEUX ALINEAS SOUS UNE SEULE CLE — t07-c1-07-1 —, et
%  ordo.py le marque « ¶2 ». On n'en ouvre pas une seconde : c'est la
%  faute que la colonne ourdoue avait faite aux tableaux 9 et 11, et
%  c'est pour elle que le marqueur existe.
%
%  LA FERME EST LE SECOND LEXIQUE ENTIEREMENT JAVANAIS DU LIVRET,
%  APRES LE CORPS ET LA PARENTE. sapi la vache (25), jaran le cheval
%  (7), wedhus le mouton et la chevre (50, 51), pitik la poule (43),
%  banyak l'oie (69), bebek le canard (66), manuk l'oiseau, kandhang
%  l'etable (58), sumur le puits (29), luku la charrue (5), garu la
%  herse (40), pacul la houe (39), arit la faucille, damen la paille,
%  suket l'herbe (80), pari le riz, gabah, lemah la terre, sawah,
%  tegal, kali la riviere (23), jugangan le sillon (8). Une langue de
%  paysans ne demande a personne le nom de ce qu'elle cultive.
%
%  ET C'EST ICI QUE « banyak » REPARAIT, CINQ TABLEAUX APRES AVOIR
%  FAIT DESARMER SA REGLE. Les oies (69) de la basse-cour se disent
%  banyak, comme celles du village du tableau 3. La regle avait ete
%  otee de kolonoj.py a ce moment-la ; sans cela le controle crierait
%  une seconde fois ici, sur une seconde forme juste. Une regle otee
%  pour une bonne raison se paie une fois et sert toujours.
%
%  DEUX \VUgras{} VIDES, ET C'EST LA FAUTE PROPRE DE LA COLONNE
%  INDONESIENNE QUI PASSE ICI. Elle s'y etait produite quatre fois
%  (tableaux 5, 9, 10 et 15), toujours a la meme place : un renvoi qui
%  suit un terme dont le groupe gras est deja ferme, la main ouvrant un
%  second groupe pour « porter » l'exposant alors que l'exposant se
%  porte tout seul. Six tableaux javanais sans une seule, puis deux
%  d'un coup au meme alinea — c'est une faute de fatigue, pas de
%  langue, et elle ne se previent pas : elle se releve. kolonoj.py l'a
%  fait des la premiere passe, sans qu'on ait rien eu a lui ajouter.
%
%  LE SEUL ENDROIT OU LE JAVANAIS EMPRUNTE DANS CE TABLEAU EST LA
%  ROUTE ET LE METIER DE VILLE : dilijens la diligence (11), pandhe le
%  forgeron (5), tukang cukur le barbier (8), asuransi l'assurance,
%  setasiun, sepatu (7), sepur. Le village fournit ses propres mots ;
%  ce qui vient de la ville arrive avec le sien. Le partage est le meme
%  qu'au tableau 6 entre le pawon et le kompor, mais tire cette fois
%  entre le dedans et le dehors du village.
%
%  LA NEIGE (37) ET LA GLACE (25) N'EXISTENT PAS A JAVA, ET LA COLONNE
%  LES ECRIT QUAND MEME. « salju » et « es » sont les mots que l'ecole
%  javanaise emploie — le premier arabe par le malais, le second
%  neerlandais (ijs) —, et ce sont les bons : la planche grave un
%  village d'Europe sous la neige. Meme raison qu'au « mangsa semi » du
%  tableau 3 : on n'echange pas une chose de la-bas contre une chose
%  d'ici sous pretexte que la langue serait plus a l'aise.
%
%  ET LE PATINEUR (26) EN EST LE CAS EXTREME. Le javanais n'a pas de
%  mot pour le patin a glace, et n'en a jamais eu besoin. La colonne
%  ecrit « sepatu es », chaussure de glace — un compose transparent,
%  fait avec deux mots empruntes, pour une chose que personne ici n'a
%  vue. C'est exactement ce que fait un manuel scolaire javanais quand
%  il enseigne l'hiver, et c'est donc la forme de 2026.
% ===================================================================

%%K t07-noto-1 noto
(*) Bab rinciné pangan, delengen gambar 6 lan 13.

%%K t07-apar-1 apar
\VUsaut{4.0mm}
\VUcentre{13.2pt}{90}{SERI KAPINDHO}
\VUsaut{3.0mm}
\VUfilet{16mm}
\VUsaut{5.6mm}
\VUtitre{15.0pt}{60}{GAMBAR No. 7}
\VUsaut{3.0mm}

%%K t07-tit-1 sub
\VUsaut{3.0mm}
\VUcentre{12.0pt}{0}{\textit{Adegan Kapisan.}}
\VUsaut{3.0mm}

%%K t07-tit-2 sub
\VUsaut{4.0mm}
\VUpk{10.2pt}{40}{Desa ing mangsa adhem.}
\VUsaut{4.0mm}

%%K t07-c1-01-1 p
1. --- Nalika prei taun anyar, aku ngeterake bapakku menyang Novurbo,
desa cilik panggonan dheweke kudu ngrampungake sawetara urusan dagang.

%%K t07-c1-02-1 p
2. --- Kita numpak \VUgras{dilijens} \textsuperscript{(11)} kang mlaku
saka kutha menyang desa iku ; nanging sarehne adhem banget, lelungan
nganggo kreta kang kurang kepenak iku ora patiya nyenengake.

%%K t07-c1-03-1 p
3. --- Mula kita rumangsa bungah temenan nalika weruh \VUgras{kusir}
mandhegake kretane ing alun-alun, ing ngarepe \VUgras{losmen}
\textsuperscript{(2)} \textit{Bruwang Putih}. \VUgras{Sing duwe losmen}
\textsuperscript{(4)} wis ngenteni kita ing tlundhake omahe, karo udud
\VUgras{pipa}ne kanthi tentrem.

%%K t07-c1-03-2 p
Dheweke ngajak kita mlebu, lan aku ora sranta lungguh ing ngarepe geni
pang garing kang krepyak-krepyak ing pawon. Ing kono aku ngguyoni
\VUgras{angin lor} kang nyocog, kang ngiyek-ngiyekake \VUgras{papan
jeneng} \textsuperscript{(3)} lawas lan nggawa lunga \VUgras{kukus}
\textsuperscript{(1)} kang metu saka cerobong kanthi ubengan ireng.

%%K t07-c1-04-1 p
4. --- Wektu iku jame sarapan. Tanpa ngenteni maneh, kita mangan enak
panganan kang prasaja nanging seger kang disuguhake marang kita (*).

%%K t07-tit-3 sub
\VUsaut{3.0mm}
\VUpk{10.2pt}{40}{Sawetara sowan.}
\VUsaut{3.0mm}

%%K t07-c1-05-1 p
5. --- Sawise sarapan, aku ngeterake bapakku, kang dadi wong asuransi,
menyang para nasabahe. Kang kapisan kita nyowani \VUgras{pandhe}
\textsuperscript{(5)}, kang manggon cedhak losmen. Dheweke lagi repot
masang ladam jaran. Aku gumun marang kekuwatane kancane kang nyekeli
kanthi kukuh \VUgras{kuku} jaran ing clemek kulite, dene pandhe masang
\VUgras{ladam} \textsuperscript{(6)} nganggo gebugan palu kang santer.

%%K t07-c1-06-1 p
6. --- Banjur kita menyang \VUgras{tukang sepatu} \textsuperscript{(7)}
desa, yaiku Iakobus kang wis tuwa. Sanajan papan jenenge nganggo
\VUgras{sepatu bot} kang endah, dheweke marem mung nyol sepatu lawas
sapasang kang wis bolong.

%%K t07-c1-06-2 p
Wong tuwa iku muni marang kita karo ngguyu : « Ing kutha, aku dadi
« tukang gawe sepatu » lan ora duwe nasabah. Ing kene, aku dadi
« tukang ndandani sepatu » lan pagaweane akeh. »

%%K t07-c1-07-1 p
7. --- Dheweke crita marang kita bab tanggane, yaiku \VUgras{tukang
cukur} \textsuperscript{(8)}, kang nasabahe padha ora marem.
\VUgras{Lading cukure} tansah sompel lan \VUgras{guntinge} ora tau
tajem.

Nanging sarehne dheweke siji-sijine tukang cukur ing desa, wong kepeksa
cukur lan ngethok rambut marang dheweke. Kejaba iku dheweke wong kang
medhit lan ora grapyak.

%%K t07-c1-08-1 p
8. --- Begjane, \VUgras{tangga} liyane ora padha karo dheweke. Tuladhane
\VUgras{tukang kreta} \textsuperscript{(9)} iku wong kang loma, sregep lan
seneng tetulung. Nyenengake ndandanake kreta marang dheweke.
Pagaweane tansah rampung.

%%K t07-c1-09-1 p
9. --- Iakobus tuwa uga ora sambat bab \VUgras{tukang pelana}
\textsuperscript{(10)}, kang kadhang dibantu njait \VUgras{abah-abah}
yen pagaweane numpuk.

%%K t07-c1-09-2 p
Bapakku takon bab kulawargane : « Kabeh padha becik. Putuku wadon ana
ing \VUgras{sekolah} \textsuperscript{(12)}. \VUgras{Guru wadone}
\textsuperscript{(13)} marem banget marang dheweke, bocahe \VUgras{murid}
\textsuperscript{(14)} kang becik. Dheweke ora padha karo anakku lanang
kang mbeling ; kang dipikir mung dolanan. \VUgras{Guru}
\textsuperscript{(15)} ora bisa nulari kawruh apa-apa marang dheweke. »

%%K t07-tit-4 sub
\VUpk{10.2pt}{40}{Gunemane wong desa.}
\VUsaut{3.0mm}

%%K t07-c1-10-1 p
10. --- Wong tuwa iku banjur crita marang kita kabar-kabare desa. Wong
bakal mbangun sekolah anyar, amarga kang ana ing \VUgras{balai desa}
wis dadi cilik banget. Kejaba iku iku gampang, amarga cukup tuku
peranganing kebone \VUgras{tukang panggilingan.} Iku bakal nguntungake
banget wong kang mesakake iku. Merga adheme, \VUgras{watu gilingan} ing
\VUgras{panggilingan} \textsuperscript{(16)} wis ora bisa nggiling
gandum, amarga \VUgras{rodhane} \textsuperscript{(17)} kandheg beku
dening \VUgras{es} \textsuperscript{(25)} ing \VUgras{kali cilik}
\textsuperscript{(23)}.

%%K t07-c1-11-1 p
11. --- « Bab bangunan, apa panjenengan wis mirsani \VUgras{greja}
\textsuperscript{(18)} anyar kita? Endah banget ing sangisore kemul
saljune. \VUgras{Manuk gagak} \textsuperscript{(19)}, kang wis ora
ngenali menara loncenge lawas, padha nyanggong kanthi sedhih ing wit
gedhe ing \VUgras{kuburan} \textsuperscript{(20)}. »

%%K t07-c1-12-1 p
12. --- Gagasan bab kuburan iku ngelingake tukang sepatu marang
wong-wong kang tepung karo bapakku lan kang mati taun kepungkur.
« Pira \VUgras{kuburan} \textsuperscript{(21)}, Pak, kang ditambahake ing
antarane liyane, ing sangisore \VUgras{wit cemara}
\textsuperscript{(22)}! Pati wis ngarit notaris kita kang lawas. Kabeh
wong desa melu \VUgras{layadan}. »

%%K t07-c1-13-1 p
13. --- « Iki penggantine kang lagi lelayaran ing kali cilik. Dheweke
lagi wae teka ndandanake tali \VUgras{sepatu es}ne
\textsuperscript{(26)} kang pedhot marang aku. »

%%K t07-c1-14-1 p
14. --- « Iki uga ing pinggire kali cilik, \VUgras{jaga desa}
\textsuperscript{(27)} anyar. Dheweke tilas prajurit polisi kang gawe
wedi bocah-bocah nakal desa. Bocah-bocah mlayu sanalika weruh
\VUgras{sarung sikile} \textsuperscript{(29)} lan \VUgras{topi kepine}
\textsuperscript{(28)}. »

%%K t07-c1-15-1 p
15. --- « Lah! iki Iosef tuwa kang nyurung \VUgras{gerobag cilik isi
kayu bongkokan} \textsuperscript{(30)} kanggo tukang roti. --- Bener
kuwi, ujare bapakku, aku ngenali jarane, yaiku \VUgras{kuldi}
\textsuperscript{(31)}. Aku pancen kudu ngomong karo dheweke, tak
gunakake kesempatane. Nganti ketemu maneh, Pak Iakobus. »

%%K t07-c1-16-1 p
16. --- Pas nalika kita metu, bocah-bocah desa mulih saka sekolahan lan
mlayu grubyug menyang \VUgras{papan luncur} \textsuperscript{(33)}
kanthi bebaya tiba lan tugel anggotane awak nalika \VUgras{tiba}
\textsuperscript{(34)}. Salah sijine njupuk \VUgras{kreta luncure}
\textsuperscript{(32)} ing tukang kreta, nglungguhake salah sijine
kancane ing kono, banjur mangkat karo mlayu.

%%K t07-c1-17-1 p
17. --- Liyane grubyug menyang \VUgras{tumpukan salju}
\textsuperscript{(37)} cedhak \VUgras{kreteg} \textsuperscript{(40)} lan
wiwit mbedhili \VUgras{reca salju} \textsuperscript{(39)} kang wingi
padha diadegake lan diwenehi topi, pipa sarta tas sekolah nyampir.

%%K t07-c1-18-1 p
18. --- Bocah-bocah ora patiya nggagas angin kang njekutake lan hawa
adhem. Kang akeh malah ora nganggo \VUgras{sarung gulu}
\textsuperscript{(35)}, lan kanggo ngangetake awak sawise perang
\VUgras{bal salju} \textsuperscript{(38)}, cukup sagegem \VUgras{kastanye}
\textsuperscript{(42)} kang panas banget kang dituku saka \VUgras{bakul}
tuwa Jakobina, kang gemeter kadhemen ing sangisore
\VUgras{slendhange} \textsuperscript{(43)} kang tipis banget.

%%K t07-tit-5 sub
\VUsaut{3.0mm}
\VUcentre{12.0pt}{0}{\textit{Adegan Kapindho.}}
\VUsaut{3.0mm}

%%K t07-tit-6 sub
\VUsaut{4.0mm}
\VUpk{10.2pt}{40}{Omah tani ing mangsa semi.}
\VUsaut{4.0mm}

%%K t07-c2-01-1 p
1. --- Sanajan ana kabeh senenge mangsa adhem, kita mapag baline
\VUgras{mangsa semi} kanthi bungah kang jero.

%%K t07-c2-01-2 p
Endah temen gegambarane latare omah tani, ing wayah sore, ing sasi Mei!

%%K t07-c2-02-1 p
2. --- Ing wates langit \VUgras{srengenge} \textsuperscript{(1)} alon-alon
angslup, ngebaki \VUgras{pedesan} \textsuperscript{(3)} nganggo
\VUgras{sunare} \textsuperscript{(2)} kang wungu. \VUgras{Wong mluku}
\textsuperscript{(6)} kang sregep gage-gage mluku \VUgras{sawahe}
\textsuperscript{(4)}.

%%K t07-c2-02-2 p
Nganggo \VUgras{luku}ne \textsuperscript{(5)} kang ditarik \VUgras{jaran}
\textsuperscript{(7)} kang kuwat, dheweke nggawe \VUgras{jugangan}
\textsuperscript{(8)}, panggonan \VUgras{wong nyebar}
\textsuperscript{(9)} bakal nyebar \VUgras{wiji} sagegem-sagegem kang
bakal ngasilake wuli emas.

%%K t07-c2-03-1 p
3. --- \VUgras{Wong ngarit} \textsuperscript{(11)} kang nyekel arit wis
ngarit sukete pangonan, tukang epe wis ngepe \VUgras{suket garing}
\textsuperscript{(10)} karo mbolak-balik nganggo \VUgras{garpu}
\textsuperscript{(12)} lan garu, lan iki wong-wong iku padha mulih.

%%K t07-c2-03-2 p
Sawetara mlaku ing ngarepe gerobag abot kang ditarik sapi ; padha
manggul pirantine ing pundhak. Liyane, kang luwih kesed, wis nglempreh
kepenak ing dhuwure suket garing kang wangi.

%%K t07-c2-04-1 p
4. --- Nalika liwat wong-wong iku aweh salam marang \VUgras{tukang
kebon} \textsuperscript{(16)} kang repot ing \VUgras{kebon woh-wohan}
\textsuperscript{(13)} ngethoki \VUgras{wit woh-wohan} lan nyambung
\VUgras{wit pir} \textsuperscript{(14)}. \VUgras{Wit plum}
\textsuperscript{(15)} wiwit kembang, lan ing \VUgras{kebon
sayur} \textsuperscript{(17)} kang wis dirabuk becik, sayuran, kubis lan
salad wis padha becik thukule.

%%K t07-c2-04-2 p
Ing tembok cendhak, \VUgras{merak} \textsuperscript{(18)} kang
endah megarake buntute supaya wulune kang endah banget dikagumi.

%%K t07-tit-7 sub
\VUsaut{3.0mm}
\VUpk{10.2pt}{40}{Latar omah tani.}
\VUsaut{3.0mm}

%%K t07-c2-05-1 p
5. --- Lawange \VUgras{kandhang jaran} \textsuperscript{(19)} padha menga
lan wong weruh papan pakan lan \VUgras{lemekan}
\textsuperscript{(23)}ne \VUgras{jaran wadon} \textsuperscript{(21)} kang
lagi wae diuculi \VUgras{tukang jaran} \textsuperscript{(20)}.
\VUgras{Belo} \textsuperscript{(22)} kang lucu temen karo sikile kang
dawa, ngetutake biyunge karo pethakilan.

%%K t07-c2-06-1 p
6. --- Cedhak \VUgras{sumur} \textsuperscript{(29)}, \VUgras{sapi wadon}
\textsuperscript{(25)} padha ngombe ing \VUgras{bak kayu}
\textsuperscript{(26)} kang dadi papan ngombene. \VUgras{Pedhet}
\textsuperscript{(24)} nganggo kesempatan iku kanggo nusu, lan
\VUgras{sapi lanang} \textsuperscript{(27)} kang ngambu wangine pakan,
mbengok bungah lan ngangkat sirahe kanthi gagah karo \VUgras{sungu}
\textsuperscript{(28)} kang kuwat.

%%K t07-c2-07-1 p
7. --- Kabeh padha nyambut gawe. \VUgras{Batur wadon}
\textsuperscript{(32)} nyekel \VUgras{tali} \textsuperscript{(31)} sumur
ing tangane. Dheweke nimba banyu nganggo \VUgras{ember}
\textsuperscript{(30)} kanggo ngombeni kewan. Cedhak \VUgras{lumbung}
\textsuperscript{(33)}, ana wong nggaruki damen, lan ing ngarepe
lawange \VUgras{omah} \textsuperscript{(34)} ana \VUgras{tukang mintal}
\textsuperscript{(35)} tuwa kang mintal \VUgras{lena} utawa
\VUgras{rami} nganggo rodha pintale lan \VUgras{antihe,} kaya ing
jaman biyen.

%%K t07-tit-8 sub
\VUsaut{3.0mm}
\VUpk{10.2pt}{40}{Wong tani.}
\VUsaut{3.0mm}

%%K t07-c2-08-1 p
8. --- \VUgras{Wong tani} \textsuperscript{(36)} kang nuntun kabeh
pagawean iku lan menehi pituduh marang para batur. Dheweke akon
nglebokake \VUgras{garu} \textsuperscript{(40)} lan \VUgras{pacul}
\textsuperscript{(39)} kang dilalekake ing cedhake \VUgras{kandhang}
\textsuperscript{(38)}ne \VUgras{asu} \textsuperscript{(37)}.

%%K t07-c2-09-1 p
9. --- Bengokane gawe kagete \VUgras{manuk dara} \textsuperscript{(41)},
kang mabur sakayange menyang \VUgras{pagupon} \textsuperscript{(42)},
lan \VUgras{pitik} \textsuperscript{(43)} kang gage-gage mlebu maneh
menyang \VUgras{kandhang pitik} \textsuperscript{(44)}.

%%K t07-c2-10-1 p
10. --- Ing ngarepe \VUgras{kandhang wedhus} \textsuperscript{(48)}, iki
\VUgras{pangon} \textsuperscript{(45)} tuwa kang nuntun bali
\VUgras{pepanthane} \textsuperscript{(49)}. Karo nyekel
\VUgras{tekene} \textsuperscript{(46)}, kang ora tau kari, kaya uga
\VUgras{klambi jaba}ne \textsuperscript{(47)} kang wis luntur, dheweke
nuntun menyang kandhang \VUgras{wedhus lanang}
\textsuperscript{(50)}, \VUgras{wedhus wadon} \textsuperscript{(53)},
\VUgras{cempe} \textsuperscript{(54)}, \VUgras{wedhus jawa wadon}
\textsuperscript{(51)} lan \VUgras{anak wedhus jawa}
\textsuperscript{(52)}. \VUgras{Asune} \textsuperscript{(55)} kang setya,
Argus, teka pungkasan lan nyurung kang keri-keri nganggo gonggongane.

%%K t07-c2-11-1 p
11. --- Kabeh rame lan obah iku ora ngganggu tentreme \VUgras{sapi
prahan} \textsuperscript{(56)} kang becik, kang ngunekake
\VUgras{kelinthinge} \textsuperscript{(57)} nalika diperes ing ngarepe
lawang \VUgras{kandhang} \textsuperscript{(58)}.

%%K t07-c2-12-1 p
12. Sapi iku kaya-kaya ngerti yen kudu menehake susune, supaya
\VUgras{bojone wong tani} \textsuperscript{(60)} bisa nyiyapake
\VUgras{mentega} ing \VUgras{pemuter mentega} \textsuperscript{(61)}
utawa \VUgras{keju} \textsuperscript{(64)} kang enak banget, kang netes
saka rak kang diselehake ing dhuwur \VUgras{bak gedhe}
\textsuperscript{(63)}.

%%K t07-c2-13-1 p
13. --- Kabeh \VUgras{papan susu} \textsuperscript{(59)} diresiki becik
banget dening \VUgras{batur tani} \textsuperscript{(65)} kang lagi wae
ngumbah \VUgras{wadhah susu} \textsuperscript{(62)} ing ember gedhe kang
digawa ing tangane.

%%K t07-tit-9 sub
\VUsaut{3.0mm}
\VUpk{10.2pt}{40}{Kandhang unggas.}
\VUsaut{3.0mm}

%%K t07-c2-14-1 p
14. --- Karo klompen kayune kang kandel, dheweke mlaku rada gedebugan,
lan mula gawe mlayune \VUgras{kalkun lanang} \textsuperscript{(68)},
\VUgras{bebek wadon} \textsuperscript{(66)}, \VUgras{bebek lanang}
\textsuperscript{(67)} lan \VUgras{banyak} \textsuperscript{(69)} kang
mangapake \VUgras{cucuke} \textsuperscript{(70)} amba lan mbengok
kanthi bingung.

%%K t07-c2-14-2 p
\VUgras{Jago} \textsuperscript{(71)}, kang luwih seneng gelut, ngadegake
\VUgras{jenggere} \textsuperscript{(72)} kang abang, ngebutake wulu
\VUgras{buntute} \textsuperscript{(73)} lan ngunekake « kukuruyuk » kang
seru banget.

%%K t07-c2-15-1 p
15. --- \VUgras{Babon} \textsuperscript{(74)} nggoleki pangan ing
\VUgras{rabuk} \textsuperscript{(81)} karo \VUgras{kuthuk-kuthuke}
\textsuperscript{(75)}. \VUgras{Gembluk} \textsuperscript{(78)} mlayu
menyang biyunge, yaiku \VUgras{babi wadon} \textsuperscript{(77)} kang
gedhe temen, kang kita deleng ing cedhake kandhang babi.

%%K t07-c2-16-1 p
16. --- Ing kandhange, \VUgras{terwelu} \textsuperscript{(79)} padha
mangan kanthi rakus \VUgras{suket} \textsuperscript{(80)} alus kang
digawakake anake wadon wong tani. Ioanina tresna banget marang
terwelune, lan saben dina, sawise njupuk \VUgras{endhog}
\textsuperscript{(76)} anyar ing kandhang pitik, dheweke teka nyowani
kanca-kanca ciliké.

\VUsaut{6.0mm}
\VUfilet{22mm}
